%-------------------------
% CV Anas MADANI — source unique, deux sorties (avec / sans photo)
%-------------------------

\documentclass[letterpaper,11pt]{article}
\usepackage{latexsym}
\usepackage[empty]{fullpage}
\usepackage{titlesec}
\usepackage{marvosym}
\usepackage[usenames,dvipsnames]{color}
\usepackage{verbatim}
\usepackage{enumitem}
\usepackage{fancyhdr}
\usepackage[utf8]{inputenc}
\usepackage[T1]{fontenc}
\usepackage{ragged2e}
\usepackage{tikz}
\usepackage{tikzfill}

% --- Liens sans cadre + metadonnees PDF (lues par certains ATS) ---
\usepackage[hidelinks]{hyperref}
\hypersetup{
  pdftitle={CV - Anas MADANI},
  pdfauthor={Anas MADANI},
  pdfsubject={Ingénieur Logiciel Full Stack},
  pdfkeywords={Java, Spring Boot, Python, FastAPI, Django, React, Docker, Kafka, PostgreSQL, Kubernetes}
}

% --- Drapeaux de compilation -------------------------------------
\newif\ifphoto
\phototrue          % <- avec photo  (Maroc)
%\photofalse        % <- sans photo  (France, remote, anglophone)

\newif\ifshowstageonea
\showstageoneatrue  % <- \showstageoneafalse pour masquer le stage 1A

\pagestyle{fancy}
\fancyhf{}
\fancyfoot{}
\renewcommand{\headrulewidth}{0pt}
\renewcommand{\footrulewidth}{0pt}

% --- Marges ------------------------------------------------------
\addtolength{\oddsidemargin}{-0.375in}
\addtolength{\evensidemargin}{-0.375in}
\addtolength{\textwidth}{1in}
\addtolength{\topmargin}{-.5in}
\addtolength{\textheight}{1.0in}

\urlstyle{same}

\raggedbottom
\raggedright
\setlength{\tabcolsep}{0in}

% --- Titres de section -------------------------------------------
\titleformat{\section}{
  \vspace{-4pt}\scshape\raggedright\large
}{}{0em}{}[\color{black}\titlerule \vspace{-5pt}]

% --- Commandes personnalisees ------------------------------------
\newcommand{\resumeItem}[2]{
  \item\small{
    \textbf{#1}{: #2 \vspace{-2pt}}
  }
}

\newcommand{\resumeSubheading}[4]{
  \vspace{-1pt}\item
    \begin{tabular*}{0.97\textwidth}{l@{\extracolsep{\fill}}r}
      \textbf{#1} & #2 \\
      \textit{\small#3} & \textit{\small #4} \\
    \end{tabular*}\vspace{-5pt}
}

\newcommand{\resumeSubItem}[2]{\resumeItem{#1}{#2}\vspace{-4pt}}

\renewcommand{\labelitemii}{$\circ$}

\newcommand{\resumeSubHeadingListStart}{\begin{itemize}[leftmargin=*]}
\newcommand{\resumeSubHeadingListEnd}{\end{itemize}}
\newcommand{\resumeItemListStart}{\begin{itemize}}
\newcommand{\resumeItemListEnd}{\end{itemize}\vspace{-5pt}}

% --- Contenu de l'en-tete, defini une seule fois ------------------
\newcommand{\cvname}{Anas MADANI}
\newcommand{\cvtitle}{Ingénieur Logiciel Full Stack}
\newcommand{\cvcontact}{
  \href{mailto:anasmadani49@gmail.com}{anasmadani49@gmail.com}
  \quad | \quad
  \href{tel:+212666322450}{+212 6 66 32 24 50}
  \quad | \quad
  Casablanca, Maroc \\[4pt]
  \href{https://www.linkedin.com/in/anas-madani/}{linkedin.com/in/anas-madani}
  \quad | \quad
  \href{https://github.com/HexNebula}{github.com/HexNebula}
  \quad | \quad 
  \href{https://hexnebula.github.io}{hexnebula.github.io}
}

%-------------------------------------------
\begin{document}

% --- En-tete -----------------------------------------------------
\ifphoto
  \begin{minipage}[c]{0.75\textwidth}
    \textbf{\Large \cvname} \\[2pt]
    \textit{\cvtitle} \\[6pt]
    \cvcontact
  \end{minipage}
  \hfill
  \begin{minipage}[c]{0.20\textwidth}
    \raggedleft
    \tikz\path[fill overzoom image={CV_light_print}]circle[radius=1.5cm];
  \end{minipage}
  \vspace{6pt}
\else
  \begin{center}
    \textbf{\Large \cvname} \\[2pt]
    \textit{\cvtitle} \\[6pt]
    \cvcontact
  \end{center}
  \vspace{4pt}
\fi

\section{Profil}
\small{
Ingénieur logiciel diplômé de l'ENSIAS (Génie Logiciel), full stack du backend au frontend, avec une pratique DevOps du déploiement en environnement Linux : intégration et livraison continues (CI/CD), conteneurisation et orchestration. Expérience en environnements exigeants (banque, datacenter), habitué à livrer en équipe agile un code testé, revu et documenté.
}

%-----------FORMATION-----------------
\section{Formation}
  \resumeSubHeadingListStart
    \resumeSubheading
      {École Nationale Supérieure d'Informatique et d'Analyse des Systèmes (ENSIAS)}{Rabat, Maroc}
      {Diplôme d'Ingénieur -- Génie Logiciel}{2022 -- 2025}
    \resumeSubheading
      {CPGE -- Settat}{Settat, Maroc}
      {Classe Préparatoire aux Grandes Écoles, TSI}{2019 -- 2022}
  \resumeSubHeadingListEnd

%-----------EXPERIENCE-----------------
\section{Expériences Professionnelles}
\resumeSubHeadingListStart

\resumeSubheading
  {Banque Centrale Populaire (BCP)}{Casablanca, Maroc}
  {Stagiaire Ingénieur Logiciel (PFE)}{Février 2025 -- Juillet 2025}
  \resumeItemListStart
    \item Conception et livraison en autonomie d'une plateforme de recommandation de produits
      bancaires : 8 microservices Spring Boot, une API de recommandation FastAPI et un
      frontend React/TypeScript, déployée en interne et présentée aux équipes métier.
    \item Développement du moteur de recommandation de bout en bout : préparation de 17 000
      paires client-produit, segmentation de clientèle par K-Means, comparaison de 4
      algorithmes de classification et d'un filtrage collaboratif par graphe (Neo4j GDS),
      avec un framework d'évaluation Precision@K / Recall@K / NDCG@K.
    \item Industrialisation de la plateforme : conteneurisation Docker, communication
      événementielle par Kafka, cache Redis, authentification Keycloak avec validation JWT
      centralisée au niveau de l'API Gateway.
    \item \textit{Stack :} Java, Spring Boot/Cloud, React, TypeScript, Python, FastAPI, Neo4j,
      PostgreSQL, Redis, Kafka, Keycloak, Docker, Git
  \resumeItemListEnd

\resumeSubheading
  {N+ONE DATACENTERS}{Casablanca, Maroc}
  {Stagiaire Ingénieur Logiciel}{Août 2024 -- Septembre 2024}
  \resumeItemListStart
    \item Développement d'une plateforme de gestion d'appels d'offres (Django REST Framework /
      React), couvrant le cycle de vie des consultations, les soumissions des candidats et un
      portail public pour les entreprises externes.
    \item Implémentation des mécanismes de sécurité : chiffrement AES-256-GCM des documents
      sensibles, contrôle d'accès par rôles, authentification JWT avec rotation des tokens,
      journal d'audit des opérations et limitation de débit sur les endpoints publics.
    \item Conteneurisation de la stack via Docker Compose (Django, PostgreSQL, React, reverse
      proxy Nginx), pipeline d'intégration continue GitHub Actions (flake8, pytest sur
      PostgreSQL) et service planifié de surveillance des échéances.
    \item \textit{Stack :} Python, Django REST Framework, React, Vite, PostgreSQL, Docker,
      Nginx, GitHub Actions, JWT, AES-GCM, pytest
  \resumeItemListEnd

\ifshowstageonea
\resumeSubheading
  {Commune Er-rich}{Er-rich, Maroc}
  {Stagiaire Ingénieur Logiciel}{Juillet 2023 -- Août 2023}
  \resumeItemListStart
    \item Développement d'une application de gestion du personnel ouvrier de la commune :
      profils des agents, suivi des absences et calcul automatisé des salaires, en
      remplacement de procédures tenues jusque-là sur papier. Seul développeur du projet.
    \item Développement full stack : services backend Spring Boot, interfaces React et
      Tailwind CSS, persistance PostgreSQL.
    \item Conception de la base de données relationnelle : dictionnaire de données,
      dépendances fonctionnelles et modèle conceptuel, à partir du recueil des besoins
      auprès du service concerné.
    \item \textit{Stack :} Java, Spring Boot, React, Tailwind CSS, PostgreSQL, UML, Git
  \resumeItemListEnd
\fi

\resumeSubHeadingListEnd

%-----------PROJETS-----------------
\section{Projets}
  \resumeSubHeadingListStart
    \resumeSubItem{Practice Debt \normalfont\textit{(Août -- Sept. 2026)}}
      {Outil de suivi d'entraînement Codeforces : identifie les problèmes abandonnés en cours de concours et les techniques non pratiquées depuis longtemps, classés dans une file unique. Rétro-ingénierie des formules de score et de classement Elo de Codeforces, erreur du modèle quantifiée contre les 20 489 changements de classement d'un seul concours Codeforces, coûts dérivés en différences modèle-contre-modèle pour annuler le biais résiduel. Backend Spring Boot 3 / PostgreSQL, frontend React 19 (TypeScript), déploiement en conteneur unique, plus de 90 tests.
      \href{https://github.com/HexNebula/practice-debt}{\underline{GitHub}}}
      \vspace{2pt}
    \resumeSubItem{BacSurv \normalfont\textit{(Août 2026)}}
      {Planification de la surveillance des examens du baccalauréat marocain : affectation des enseignants aux séances sous contraintes de disponibilité, de charge et de conflits. Troisième version d'un problème abordé pour la première fois en 2025 -- un prototype TypeScript, puis une réécriture FastAPI, aujourd'hui Java/Spring Boot avec solveur de contraintes. Frontend React 19 (TypeScript), export Word/Excel et interface bilingue arabe/français.
      \href{https://github.com/HexNebula/BacSurv}{\underline{GitHub}}}
  
      \vspace{2pt}
    \resumeSubItem{BusConnect \normalfont\textit{(Oct. 2024 -- Janv. 2025)}}
      {Plateforme de réservation de bus en architecture microservices avec Spring Boot 3, Spring Cloud (Eureka, API Gateway), React 18 et déploiement Kubernetes. Intègre PostgreSQL, MySQL et MongoDB selon les services.
      \href{https://github.com/HexNebula/BusManagementV2}{\underline{GitHub}}}
      \vspace{2pt}
    \resumeSubItem{MLOps Pipeline Generator \normalfont\textit{(Oct. 2024 -- Janv. 2025)}}
      {Générateur de code MLOps basé sur l'ingénierie dirigée par les modèles (MDE) avec Eclipse Ecore et Acceleo. Génère automatiquement des scripts Python, Dockerfiles et configurations Kubernetes à partir de métamodèles via une interface SWT/EMF Forms.
      \href{https://github.com/HexNebula/MLOps-Pipeline-Code-Generator}{\underline{GitHub}}}
\resumeSubHeadingListEnd

%-----------COMPETENCES-----------------
\section{Compétences}
  \resumeSubHeadingListStart
    \item{
      \textbf{Langages}{: Java, Python, TypeScript, JavaScript, C, C++, SQL} \\
      \textbf{Backend}{: Spring Boot, Spring Cloud, FastAPI, Django REST Framework} \\
      \textbf{Frontend}{: React, TypeScript, Tailwind CSS, Vite} \\
      \textbf{Bases de données}{: PostgreSQL, MySQL, MongoDB, Neo4j, Redis} \\
      \textbf{DevOps \& Infra}{: Docker, Kubernetes, Kafka, GitHub Actions, Nginx, Keycloak, Linux} \\
      \textbf{Langues}{: Arabe (natif), Français (courant), Anglais (courant)}
    }
  \resumeSubHeadingListEnd

\end{document}
